% Options for packages loaded elsewhere
% Options for packages loaded elsewhere
\PassOptionsToPackage{unicode}{hyperref}
\PassOptionsToPackage{hyphens}{url}
\PassOptionsToPackage{dvipsnames,svgnames,x11names}{xcolor}
%
\documentclass[
  spanish,
  10pt,
  letterpaper,
  twocolumn]{article}
\usepackage{xcolor}
\usepackage[top=1.0cm,bottom=1.5cm,left=1.5cm,right=1.5cm]{geometry}
\usepackage{amsmath,amssymb}
\setcounter{secnumdepth}{-\maxdimen} % remove section numbering
\usepackage{iftex}
\ifPDFTeX
  \usepackage[T1]{fontenc}
  \usepackage[utf8]{inputenc}
  \usepackage{textcomp} % provide euro and other symbols
\else % if luatex or xetex
  \usepackage{unicode-math} % this also loads fontspec
  \defaultfontfeatures{Scale=MatchLowercase}
  \defaultfontfeatures[\rmfamily]{Ligatures=TeX,Scale=1}
\fi
\usepackage{lmodern}
\ifPDFTeX\else
  % xetex/luatex font selection
  \setmainfont[]{Fira Sans}
\fi
% Use upquote if available, for straight quotes in verbatim environments
\IfFileExists{upquote.sty}{\usepackage{upquote}}{}
\IfFileExists{microtype.sty}{% use microtype if available
  \usepackage[]{microtype}
  \UseMicrotypeSet[protrusion]{basicmath} % disable protrusion for tt fonts
}{}
\usepackage{setspace}
\makeatletter
\@ifundefined{KOMAClassName}{% if non-KOMA class
  \IfFileExists{parskip.sty}{%
    \usepackage{parskip}
  }{% else
    \setlength{\parindent}{0pt}
    \setlength{\parskip}{6pt plus 2pt minus 1pt}}
}{% if KOMA class
  \KOMAoptions{parskip=half}}
\makeatother
% Make \paragraph and \subparagraph free-standing
\makeatletter
\ifx\paragraph\undefined\else
  \let\oldparagraph\paragraph
  \renewcommand{\paragraph}{
    \@ifstar
      \xxxParagraphStar
      \xxxParagraphNoStar
  }
  \newcommand{\xxxParagraphStar}[1]{\oldparagraph*{#1}\mbox{}}
  \newcommand{\xxxParagraphNoStar}[1]{\oldparagraph{#1}\mbox{}}
\fi
\ifx\subparagraph\undefined\else
  \let\oldsubparagraph\subparagraph
  \renewcommand{\subparagraph}{
    \@ifstar
      \xxxSubParagraphStar
      \xxxSubParagraphNoStar
  }
  \newcommand{\xxxSubParagraphStar}[1]{\oldsubparagraph*{#1}\mbox{}}
  \newcommand{\xxxSubParagraphNoStar}[1]{\oldsubparagraph{#1}\mbox{}}
\fi
\makeatother


\usepackage{longtable,booktabs,array}
\newcounter{none} % for unnumbered tables
\usepackage{calc} % for calculating minipage widths
% Correct order of tables after \paragraph or \subparagraph
\usepackage{etoolbox}
\makeatletter
\patchcmd\longtable{\par}{\if@noskipsec\mbox{}\fi\par}{}{}
\makeatother
% Allow footnotes in longtable head/foot
\IfFileExists{footnotehyper.sty}{\usepackage{footnotehyper}}{\usepackage{footnote}}
\makesavenoteenv{longtable}
\usepackage{graphicx}
\makeatletter
\newsavebox\pandoc@box
\newcommand*\pandocbounded[1]{% scales image to fit in text height/width
  \sbox\pandoc@box{#1}%
  \Gscale@div\@tempa{\textheight}{\dimexpr\ht\pandoc@box+\dp\pandoc@box\relax}%
  \Gscale@div\@tempb{\linewidth}{\wd\pandoc@box}%
  \ifdim\@tempb\p@<\@tempa\p@\let\@tempa\@tempb\fi% select the smaller of both
  \ifdim\@tempa\p@<\p@\scalebox{\@tempa}{\usebox\pandoc@box}%
  \else\usebox{\pandoc@box}%
  \fi%
}
% Set default figure placement to htbp
\def\fps@figure{htbp}
\makeatother



\ifLuaTeX
\usepackage[bidi=basic,shorthands=off]{babel}
\else
\usepackage[bidi=default,shorthands=off]{babel}
\fi
\ifPDFTeX
\else
\babelfont{rm}[]{Fira Sans}
\fi
\ifLuaTeX
  \usepackage{selnolig} % disable illegal ligatures
\fi


\setlength{\emergencystretch}{3em} % prevent overfull lines

\providecommand{\tightlist}{%
  \setlength{\itemsep}{0pt}\setlength{\parskip}{0pt}}



 


\usepackage{tikz}
\usepackage{graphicx}
\usepackage{amsmath}
\usepackage{multirow}
\renewcommand{\arraystretch}{1.4}
\makeatletter
\@ifpackageloaded{caption}{}{\usepackage{caption}}
\AtBeginDocument{%
\ifdefined\contentsname
  \renewcommand*\contentsname{Tabla de contenidos}
\else
  \newcommand\contentsname{Tabla de contenidos}
\fi
\ifdefined\listfigurename
  \renewcommand*\listfigurename{Listado de Figuras}
\else
  \newcommand\listfigurename{Listado de Figuras}
\fi
\ifdefined\listtablename
  \renewcommand*\listtablename{Listado de Tablas}
\else
  \newcommand\listtablename{Listado de Tablas}
\fi
\ifdefined\figurename
  \renewcommand*\figurename{Figura}
\else
  \newcommand\figurename{Figura}
\fi
\ifdefined\tablename
  \renewcommand*\tablename{Tabla}
\else
  \newcommand\tablename{Tabla}
\fi
}
\@ifpackageloaded{float}{}{\usepackage{float}}
\floatstyle{ruled}
\@ifundefined{c@chapter}{\newfloat{codelisting}{h}{lop}}{\newfloat{codelisting}{h}{lop}[chapter]}
\floatname{codelisting}{Listado}
\newcommand*\listoflistings{\listof{codelisting}{Listado de Listados}}
\makeatother
\makeatletter
\makeatother
\makeatletter
\@ifpackageloaded{caption}{}{\usepackage{caption}}
\@ifpackageloaded{subcaption}{}{\usepackage{subcaption}}
\makeatother
\usepackage{bookmark}
\IfFileExists{xurl.sty}{\usepackage{xurl}}{} % add URL line breaks if available
\urlstyle{same}
% fallback for those not using the hyperref driver hyperxmp:
\makeatletter
\@ifundefined{xmpquote}{\newcommand{\xmpquote}[1]{#1}}{}
\makeatother
\hypersetup{
  pdftitle={Práctica 7 -- Fisiología de las Fitohormonas},
  pdflang={es},
  colorlinks=true,
  linkcolor={blue},
  filecolor={Maroon},
  citecolor={Blue},
  urlcolor={Blue},
  pdfcreator={LaTeX via pandoc}}


\title{\vspace{-1.5em}{\Huge \textbf{Práctica 7 -- Fisiología de las
Fitohormonas}}\vspace{-1.5em}}
\date{}
\begin{document}
\maketitle


\setstretch{1.2}
\textbf{Fecha}: \hrulefill \hfill \textbf{Grupo}: \hrulefill

\textbf{Nombre:} \hrulefill

\textbf{Objetivos:}

\begin{itemize}
\tightlist
\item
  Evaluar el efecto biofisiológico del etileno emitido por una fuente
  natural (manzana roja madura) sobre la germinación y el crecimiento
  inicial de plántulas de lenteja o frijol.
\item
  Identificar la caracterización morfológica de la \textbf{``triple
  respuesta''} al etileno en dicotiledóneas (inhibición de la elongación
  caulinar, engrosamiento del hipocótilo y crecimiento
  horizontal/exageración del gancho apical).
\end{itemize}

\subsection{1. Materiales}\label{materiales}

\begin{itemize}
\tightlist
\item
  60 semillas de lenteja (\emph{Lens culinaris}) o frijol
  (\emph{Phaseolus vulgaris}).
\item
  6 vasos plásticos o macetas pequeñas con sustrato húmedo (algodón o
  vermiculita).
\item
  2 manzanas rojas maduras.
\item
  2 bolsas plásticas negras de basura (grandes y herméticas).
\item
  Regla graduada milimetrada y pie de rey / calibrador.
\item
  Marcador permanente y cinta de enmascarar.
\item
  Agua destilada o corriente.
\end{itemize}

\subsection{2. Aspectos Teóricos e
Introducción}\label{aspectos-teuxf3ricos-e-introducciuxf3n}

El \textbf{etileno} (\(C_2H_4\)) es una fitohormona gaseosa única que
regula procesos como la maduración de frutos climatéricos, la
senescencia, la abscisión foliar y las respuestas a estrés mecánico. En
plántulas etioladas de dicotiledóneas, la acumulación de etileno
desencadena la clásica \textbf{triple respuesta}, una adaptación
morfológica que permite a la plántula esquivar obstáculos subterráneos
durante la emergencia.

Defina brevemente:

\begin{enumerate}
\def\labelenumi{\arabic{enumi}.}
\tightlist
\item
  \textbf{Mecanismo Fisiológico y Síntomas de la Triple Respuesta al
  Etileno:}
\end{enumerate}

\hrulefill

\hrulefill

\hrulefill

\hrulefill

\hrulefill

\hrulefill

\begin{enumerate}
\def\labelenumi{\arabic{enumi}.}
\setcounter{enumi}{1}
\tightlist
\item
  \textbf{Frutos Climatéricos vs.~No Climatéricos y Emisión de Etileno:}
\end{enumerate}

\hrulefill

\hrulefill

\hrulefill

\hrulefill

\hrulefill

\hrulefill

\begin{enumerate}
\def\labelenumi{\arabic{enumi}.}
\setcounter{enumi}{2}
\tightlist
\item
  \textbf{Evolución del Gancho Apical y Protección del Meristemo:}
\end{enumerate}

\hrulefill

\hrulefill

\hrulefill

\hrulefill

\hrulefill

\hrulefill

\subsection{3. Procedimiento
Experimental}\label{procedimiento-experimental}

\subsubsection{A. Siembra e Imbibición de
Semillas}\label{a.-siembra-e-imbibiciuxf3n-de-semillas}

\begin{enumerate}
\def\labelenumi{\arabic{enumi}.}
\tightlist
\item
  Seleccione semillas homogéneas de lenteja o frijol y distribúyalas en
  6 vasos o macetas pequeñas (10 semillas por maceta) sobre sustrato
  humedecido uniformemente.
\end{enumerate}

\subsubsection{B. Montaje de los Microinvernaderos
Oscuros}\label{b.-montaje-de-los-microinvernaderos-oscuros}

\begin{enumerate}
\def\labelenumi{\arabic{enumi}.}
\tightlist
\item
  \textbf{Tratamiento 1 (Control - Sin Etileno):} Coloque 3 macetas
  (Réplicas A, B, C) dentro de la primera bolsa negra de basura. Selle
  la bolsa herméticamente dejando aire retenido en su interior.
\item
  \textbf{Tratamiento 2 (Etileno - Con Manzana Madura):} Coloque las 3
  macetas restantes (Réplicas A, B, C) dentro de la segunda bolsa negra
  de basura e introduzca 2 manzanas rojas maduras al lado de los vasos.
  Selle la bolsa herméticamente reteniendo aire.
\item
  Mantenga ambos sistemas en un lugar con temperatura constante durante
  7 días en oscuridad ininterrumpida.
\end{enumerate}

\begin{center}
\begin{tikzpicture}[scale=0.85]
  % TRATAMIENTO 1: Control (Sin Manzana)
  \begin{scope}[shift={(0,0)}]
    % Bolsa negra contenedor
    \draw[thick, fill=gray!85, rounded corners=8pt] (-0.2,-0.2) rectangle (2.2, 2.3);
    \node[font=\fontsize{5}{6}\bfseries\selectfont, white] at (1.0, 2.1) {Control};
    
    % Macetas (n=3)
    \foreach \x in {0.1, 0.75, 1.4} {
      \draw[thick, fill=brown!40] (\x,0.1) -- (\x+0.1, 0.7) -- (\x+0.5, 0.7) -- (\x+0.6, 0.1) -- cycle;
      % Plántulas etioladas rectas
      \draw[line width=1.2pt, green!60!black] (\x+0.3,0.7) -- (\x+0.3, 1.6);
      \fill[green!50!black] (\x+0.28,1.6) circle (2pt);
    }

    \node[font=\fontsize{6}{7}\bfseries\selectfont, align=center] at (1.0, -0.4) {T1: Control};
    \node[font=\fontsize{5}{6}\selectfont, gray!50!black, align=center] at (1.0, -0.7) {Sin Manzana ($n=3$)};
  \end{scope}

  % TRATAMIENTO 2: Etileno (Bolsa Negra con Manzana)
  \begin{scope}[shift={(3.4,0)}]
    % Bolsa negra contenedor
    \draw[thick, fill=gray!85, rounded corners=8pt] (-0.2,-0.2) rectangle (2.6, 2.3);
    \node[font=\fontsize{5}{6}\bfseries\selectfont, white] at (1.2, 2.1) {Etileno};
    
    % Macetas (n=3)
    \foreach \x in {0.0, 0.65, 1.3} {
      \draw[thick, fill=brown!40] (\x,0.1) -- (\x+0.1, 0.7) -- (\x+0.5, 0.7) -- (\x+0.6, 0.1) -- cycle;
      % Plántulas con Triple Respuesta (cortas, gruesas y con gancho)
      \draw[line width=2.2pt, green!50!black] (\x+0.3,0.7) .. controls (\x+0.3, 1.1) .. (\x+0.15, 1.15);
      \fill[green!50!black] (\x+0.15,1.15) circle (2.2pt);
    }
    
    % Manzana roja (Fuente de etileno)
    \fill[red!80!black] (2.15, 0.4) circle (6pt);
    \draw[green!30!black, line width=1pt] (2.15, 0.6) -- (2.15, 0.7);
    \node[font=\fontsize{4}{5}\selectfont, white] at (2.15, 0.38) {Etileno};

    \node[font=\fontsize{6}{7}\bfseries\selectfont, align=center] at (1.2, -0.4) {T2: Etileno};
    \node[font=\fontsize{5}{6}\selectfont, gray!50!black, align=center] at (1.2, -0.7) {Con Manzana Roja ($n=3$)};
  \end{scope}
\end{tikzpicture}
\end{center}

\subsubsection{C. Evaluación Cuantitativa (Día
7)}\label{c.-evaluaciuxf3n-cuantitativa-duxeda-7}

\begin{enumerate}
\def\labelenumi{\arabic{enumi}.}
\tightlist
\item
  Abra con cuidado las bolsas al cabo de 7 días.
\item
  Registre el \textbf{Porcentaje de Germinación} (\(\%G\)) en cada
  maceta.
\item
  En 5 plántulas seleccionadas al azar por maceta, mida con la regla la
  \textbf{longitud del vástago/hipocótilo} (\(\text{cm}\)) y con el pie
  de rey el \textbf{grosor del hipocótilo} (\(\text{mm}\)).
\end{enumerate}

\subsection{4. Registro de Resultados}\label{registro-de-resultados}

En la siguiente tabla registre los resultados observados al séptimo día
del porcentaje de germinación (Germ. \(\%\)), longitud del vástago (L.
Vast. mm) y grosor del hipocótilo (Gros. Hip. mm).

Calcule el promedio (\(\bar{x}\)) y la desviación estándar (\(s\)) para
cada variable de respuesta:

\[\bar{x} = \frac{\sum_{i=1}^{n} x_i}{n} \qquad s = \sqrt{\frac{\sum_{i=1}^{n} (x_i - \bar{x})^2}{n - 1}}\]

\subsubsection{Tabla 1. Variables Morfométricas de Germinación y
Desarrollo al Día
7}\label{tabla-1.-variables-morfomuxe9tricas-de-germinaciuxf3n-y-desarrollo-al-duxeda-7}

\begin{center}
\resizebox{\linewidth}{!}{%
\begin{tabular}{p{2.3cm}c|c|c|c}
\hline
\multicolumn{1}{p{2.3cm}|}{\textbf{Tratamiento}} & \textbf{Rép.} & \textbf{Germ. ($\%$)} & \textbf{L. Vást. (mm)} & \textbf{Gros. Hip. (mm)} \\ \hline
\multicolumn{1}{p{2.3cm}|}{\multirow{3}{=}{1. Control\newline (Sin manzana)}} & A &  &  &  \\
\multicolumn{1}{p{2.3cm}|}{} & B &  &  &  \\
\multicolumn{1}{p{2.3cm}|}{} & C &  &  &  \\ \hline
\multicolumn{2}{c|}{\textbf{Promedio ($\bar{x}$)}} &  &  &  \\
\multicolumn{2}{c|}{\textbf{Desv. Est. ($s$)}} &  &  &  \\ \hline\hline
\multicolumn{1}{p{2.3cm}|}{\multirow{3}{=}{2. Etileno\newline (Con manzana)}} & A &  &  &  \\
\multicolumn{1}{p{2.3cm}|}{} & B &  &  &  \\
\multicolumn{1}{p{2.3cm}|}{} & C &  &  &  \\ \hline
\multicolumn{2}{c|}{\textbf{Promedio ($\bar{x}$)}} &  &  & \\
\multicolumn{2}{c|}{\textbf{Desv. Est. ($s$)}} &  &  & \\ \hline
\end{tabular}%
}
\end{center}

Interprete los resultados del porcentaje de germinación, longitud del
vástago y grosor del hipocótilo:

\hrulefill

\hrulefill

\hrulefill

\hrulefill

\hrulefill

\hrulefill

\subsubsection{Tabla 2. Observaciones Cualitativas del Desarrollo
Morfológico}\label{tabla-2.-observaciones-cualitativas-del-desarrollo-morfoluxf3gico}

\begin{center}
\resizebox{\linewidth}{!}{%
\begin{tabular}{l|p{14em}}
\hline
\textbf{Tratamiento} & \textbf{Curvatura del Gancho Apical, Orientación de Crecimiento y Coloración} \\
\hline
\shortstack[l]{\\1. Control\\(Sin manzana)}& \\[10ex]
\hline
\shortstack[l]{\\2. Etileno\\(Con Manzana)} & \\[10ex]
\hline
\end{tabular}%
}
\end{center}

Interprete los resultados de las observaciones cualitativas del
desarrollo morfológico:

\hrulefill

\hrulefill

\hrulefill

\hrulefill

\hrulefill

\hrulefill

\subsection{5. Análisis Gráfico}\label{anuxe1lisis-gruxe1fico}

Represente la longitud del vástago y el grosor del hipocótilo utilizando
diagramas de cajas y bigotes (boxplots) por tratamiento.

\subsubsection{Gráfica 1: Longitud del Vástago
(mm)}\label{gruxe1fica-1-longitud-del-vuxe1stago-mm}

\begin{center}
\begin{tikzpicture}[scale=1.05]
  % Grid ggplot2
  \draw[gray!20, step=0.7cm] (0,0) grid (7,5);
  \draw[gray!50, step=1.4cm] (0,0) grid (7,5);
  \draw[thick] (0,0) rectangle (7,5);
  
  % Eje X
  \node[below, font=\small] at (2.1,0) {\textbf{Control}};
  \node[below, font=\small] at (4.9,0) {\textbf{Etileno}};
  \node[below=11pt, font=\small] at (3.25,0) {\textbf{Tratamiento Experimental}};
  
  % Eje Y
  \node[rotate=90, above=10pt, font=\small] at (0,2.5) {\textbf{Longitud del Vástago (mm)}};
\end{tikzpicture}
\end{center}

\subsubsection{Gráfica 2: Grosor del Hipocótilo
(mm)}\label{gruxe1fica-2-grosor-del-hipocuxf3tilo-mm}

\begin{center}
\begin{tikzpicture}[scale=1.05]
  % Grid ggplot2
  \draw[gray!20, step=0.7cm] (0,0) grid (7,5);
  \draw[gray!50, step=1.4cm] (0,0) grid (7,5);
  \draw[thick] (0,0) rectangle (7,5);
  
  % Eje X
  \node[below, font=\small] at (2.1,0) {\textbf{Control}};
  \node[below, font=\small] at (4.9,0) {\textbf{Etileno}};
  \node[below=11pt, font=\small] at (3.25,0) {\textbf{Tratamiento Experimental}};
  
  % Eje Y
  \node[rotate=90, above=10pt, font=\small] at (0,2.5) {\textbf{Grosor del Hipocótilo (mm)}};
\end{tikzpicture}
\end{center}

Interprete los resultados representados en las gráficas:

\hrulefill

\hrulefill

\hrulefill

\hrulefill

\hrulefill

\hrulefill

\subsection{6. Actividades
Complementarias}\label{actividades-complementarias}

\begin{enumerate}
\def\labelenumi{\arabic{enumi}.}
\tightlist
\item
  Explique el papel bioquímico del precursor
  1-aminociclopropano-1-carboxílico (\textbf{ACC}) y de la \textbf{ACC
  oxidasa} en la vía biosintética del etileno en las manzanas.
\end{enumerate}

\hrulefill

\hrulefill

\hrulefill

\hrulefill

\hrulefill

\hrulefill

\begin{enumerate}
\def\labelenumi{\arabic{enumi}.}
\setcounter{enumi}{1}
\tightlist
\item
  ¿Cuál es el significado ecológico adaptativo del engrosamiento del
  hipocótilo y la curvatura del gancho apical cuando una plántula
  germina bajo un suelo compactado?
\end{enumerate}

\hrulefill

\hrulefill

\hrulefill

\hrulefill

\hrulefill

\hrulefill

\begin{enumerate}
\def\labelenumi{\arabic{enumi}.}
\setcounter{enumi}{2}
\tightlist
\item
  ¿Por qué es fundamental emplear bolsas negras herméticas en este
  experimento tanto para el control como para el tratamiento con
  etileno?
\end{enumerate}

\hrulefill

\hrulefill

\hrulefill

\hrulefill

\hrulefill

\hrulefill

\subsection{7. Conclusión}\label{conclusiuxf3n}

Compare los promedios obtenidos entre tratamientos e interprete la
presencia o ausencia de la triple respuesta al etileno en las plántulas:

\hrulefill

\hrulefill

\hrulefill

\hrulefill

\hrulefill

\subsection{8. Bibliografía}\label{bibliografuxeda}

\begin{itemize}
\tightlist
\item
  Abeles, F. B., Morgan, P. W., \& Saltveit, M. E. (1992).
  \emph{Ethylene in Plant Biology} (2nd ed.). Academic Press.
\item
  Azcón-Bieto, J., \& Talón, M. (2008). \emph{Fundamentos de Fisiología
  Vegetal} (2.ª ed.). McGraw-Hill.
\item
  Taiz, L., \& Zeiger, E. (2002). \emph{Plant Physiology} (3rd ed.).
  Sinauer Associates.
\end{itemize}




\end{document}
